% =====================================================================
\chapter{Autoevaluación}
\label{ap:autoevaluacion}
% =====================================================================

\section{Puntuación por parte y por módulo}

El examen se corrige con la clave del \cref{tab:clave} y con los criterios entre corchetes de cada solución. El \cref{tab:puntos-parte} recoge la puntuación por parte y el \cref{tab:puntos-modulo}, la puntuación por módulo, que sirve para localizar las debilidades: cada módulo vale 20 puntos.

\begin{table}[h]
\centering\small
\caption{Puntuación por parte.}\label{tab:puntos-parte}
\begin{tabular}{@{}lcc@{}}
\toprule
Parte & Puntos & Nota obtenida \\
\midrule
A. Cuestiones (20 de 2 puntos) & 40 & \\
B. Problemas (5 de 12 puntos) & 60 & \\
\midrule
\textbf{Examen} & \textbf{100} & \\
\midrule
C. Proyecto final (se califica aparte) & 100 & \\
\bottomrule
\end{tabular}
\end{table}

\begin{table}[h]
\centering\small
\caption{Puntuación por módulo.}\label{tab:puntos-modulo}
\begin{tabularx}{\textwidth}{@{}>{\raggedright\arraybackslash}X>{\raggedright\arraybackslash}p{2.6cm}>{\raggedright\arraybackslash}p{2.2cm}ccc@{}}
\toprule
Módulo & Cuestiones & Problema & Puntos & Umbral 70\,\% & Obtenidos \\
\midrule
M24. Formulación y resolución & 1--4 (8~p) & B.1 (12~p) & 20 & 14 & \\
M25. Optimización bayesiana y sustitutos & 5--8 (8~p) & B.2 (12~p) & 20 & 14 & \\
M26. Control predictivo y RL & 9--12 (8~p) & B.3 (12~p) & 20 & 14 & \\
M27. Geometría y 3D & 13--16 (8~p) & B.4 (12~p) & 20 & 14 & \\
M28. Búsqueda con LLM, producción y valor & 17--20 (8~p) & B.5 (12~p) & 20 & 14 & \\
\midrule
\textbf{Total} & 20 cuestiones (40~p) & 5 problemas (60~p) & \textbf{100} & \textbf{70} & \\
\bottomrule
\end{tabularx}
\end{table}

El \cref{eje:simp} usa en el apartado~(e) las condiciones de Karush--Kuhn--Tucker de M24, §4.2; se asigna entero a M27 porque el resto del problema es optimización de topología.

\section{Umbral de aprobado}

\begin{itemize}[label=\textendash]
  \item \textbf{Examen de hito}: 70 puntos de 100 (70\,\%) o más.
  \item \textbf{Proyecto final}: 70 puntos de 100 o más en la rúbrica de la opción elegida y todos los criterios de aceptación obligatorios cumplidos (\cref{sec:proy-calificacion}).
  \item \textbf{Puerta de salida de la fase~VI}: además del examen, cada uno de los módulos M24 a M28 exige un test de módulo con al menos 9 respuestas correctas de 12 y una autoevaluación de las prácticas o del laboratorio interactivo de la fase de al menos el 70\,\%.
  \item \textbf{Nota del hito}: $0{,}4\cdot\text{examen}/10+0{,}6\cdot\text{proyecto}/10$, con al menos un 7 en cada parte.
\end{itemize}

Como orientación por módulo, una puntuación inferior a 14 de 20 en el \cref{tab:puntos-modulo} indica que conviene repasar ese módulo aunque el examen esté aprobado. Como la fase~VI cierra el curso, superar este hito completa «Anatomía de un LLM».

\section{Qué hacer si no se alcanza el umbral}

\begin{enumerate}[label=\arabic*.]
  \item Anotar cada fallo en el registro de errores clasificándolo como error de concepto, de cálculo o despiste.
  \item Repasar las secciones del \cref{tab:repaso} correspondientes a los fallos, empezando por los errores de concepto.
  \item Rehacer sin mirar la solución los ejercicios de nivel $\star\star$ y $\star\star\star$ de esos temas en los apuntes del módulo y repetir el caso de estudio correspondiente en su cuaderno de prácticas.
  \item Repetir el test del módulo afectado hasta alcanzar 9 de 12 y volver a presentarse al examen de hito con una versión nueva, no antes de una semana.
  \item Si falla el proyecto, corregir los criterios obligatorios incumplidos con las orientaciones del \cref{ap:correccion} y volver a entregar en un plazo de cuatro semanas; se conserva la nota de los criterios ya superados.
\end{enumerate}

\begin{table}[h]
\centering\small
\caption{Secciones de repaso según el fallo (comprobadas en los índices de los apuntes).}\label{tab:repaso}
\begin{tabularx}{\textwidth}{@{}>{\raggedright\arraybackslash}p{4.2cm}X@{}}
\toprule
Fallo & Repasar \\
\midrule
Cuestiones 1, 2; problema B.1 & M24, §2.1 (programas lineales), §2.2 (dualidad y precios sombra), §2.3 (programación entera mixta, big-M) y §2.4 (ramificación y acotación) \\
Cuestión 3 & M24, §6.2 (vendedor de periódicos) y §6.1 (predecir y después optimizar) \\
Cuestión 4 & M24, §5.4 (optimización multiobjetivo) \\
Cuestiones 5, 8; problema B.2 (c)--(d) & M25, §4.2 (probabilidad de mejora), §4.3 (mejora esperada), §4.4 (cotas de confianza) y §4.5 (varianza latente y predictiva) \\
Cuestión 6; problema B.2 (a), (b), (e) & M25, §3.3 (procesos gaussianos) y §3.4 (hiperparámetros) \\
Cuestión 7 & M25, §6.2 (búsqueda aleatoria frente a rejilla) \\
Cuestión 9 & M26, §1.2 (discretización) y §1.3 (estabilidad) \\
Cuestión 10; problema B.3 (a) & M26, §3.1 (regulador lineal cuadrático) \\
Cuestión 11; problema B.3 (b) & M26, §3.2 (horizonte deslizante) y §3.3 (MPC como programa cuadrático, penalización exacta) \\
Cuestión 12; problema B.3 (c)--(d) & M26, §4.1 (ecuación de Bellman) y §4.4 (Gymnasium: \texttt{terminated} y \texttt{truncated}); M21, §1.3 \\
Cuestiones 13, 14; problema B.4 & M27, §2.1 (flexibilidad mínima), §2.2 (SIMP), §2.3 (método adjunto) y §2.4 (criterios de optimalidad); M24, §4.2 (KKT) \\
Cuestión 15 & M27, §5.3 (error en pantalla) y §5.4 (selección con presupuesto) \\
Cuestión 16 & M27, §1.2 (mallas y característica de Euler) \\
Cuestión 17 & M28, §2.2 (sesgo del ganador) \\
Cuestión 18 & M28, §4.3 (deriva y reentrenamiento) \\
Cuestión 19 & M28, §3.3 (multifidelidad y coste) \\
Cuestión 20; problema B.5 & M28, §5.1 (coste y beneficio), §5.2 (intervalos para la mejora) y §5.3 (caso de estudio); 0.9, §4.4 y §5.4 \\
\bottomrule
\end{tabularx}
\end{table}

\paragraph{Siguiente paso.} Superado el hito, el curso está completo. El proyecto final, con su documento de diseño, su código y su informe, es la pieza que conviene conservar como muestra de trabajo.
